\documentclass[11pt]{article}
% Préambule commun des prétirages v2 (format arXiv : source TeX autonome, pdfLaTeX).
% Moteur : pdfLaTeX (arXiv) → fontenc T1 + inputenc ; XeTeX/LuaTeX (tectonic) → fontspec. Sans ce choix, les
% guillemets « » tapés en dur sortent en « ń ż » sous tectonic (ledger, 2026-11-14).
\usepackage{iftex}
\ifPDFTeX
  \usepackage[T1]{fontenc}
  \usepackage[utf8]{inputenc}
\else
  \usepackage{fontspec}
\fi
\usepackage[french]{babel}
\usepackage{lmodern,microtype}
\usepackage[a4paper,margin=2.5cm]{geometry}
\usepackage{amsmath,amssymb}
\usepackage{graphicx,booktabs,siunitx}
\usepackage{authblk}
\usepackage[numbers,sort&compress]{natbib}
\usepackage[hidelinks]{hyperref}
\usepackage{csquotes}
\sisetup{locale=FR,group-minimum-digits=4}
\renewcommand\Affilfont{\small}
% Étiquette de provenance obligatoire sur chaque nombre (règle du projet) :
\newcommand{\mesure}{\textsuperscript{\textsc{m}}}   % mesuré
\newcommand{\derive}{\textsuperscript{\textsc{d}}}   % dérivé
\newcommand{\suppose}{\textsuperscript{\textsc{s}}}  % supposé


\title{Prévoir l'effectif libéral des professions de santé par département :\\
un shrinkage hiérarchique améliore de peu le stock-flux, les âges fins non}
\author[1]{Équipe hackathon-santé}
\author[1]{Claude (Anthropic), assistant de recherche}
\affil[1]{Projet « jumeau numérique de l'offre de soins », dépôt \texttt{hackathon-sante/v2}}
\date{Prétirage, version du \today{} --- non relu par des pairs}

\begin{document}
\sloppy
\maketitle

\begin{abstract}
On cherche une méthode interprétable qui prévoie, mieux que le modèle stock-flux à taux locaux,
l'effectif \enquote{libéral} (secteur privé, lucratif ou non) des sages-femmes, chirurgiens-dentistes,
pharmaciens, masseurs-kinésithérapeutes et infirmiers par département. Cinq variantes ont été
inscrites avant tout passage sur un banc de prévision pré-enregistré (origines glissantes 2016--2025,
erreur absolue moyenne en pourcentage, MAPE, par département à cinq ans). Les cinq sont des stock-flux :
même équation, estimation différente. Un shrinkage hiérarchique des taux (département, région,
national) donne une MAPE de 8,57\,\%\mesure{} contre 8,68\,\%\mesure{} pour la référence, soit
$-0{,}11$\mesure{} point (IC 95\,\% $[-0{,}21\,;\,-0{,}02]$\mesure), et le gain se reproduit sur un second
périmètre ($-0{,}10$\mesure{} point). Des âges en années simples font jeu égal sur l'effectif total ;
sur des mondes simulés, ils divisent par trois\mesure{} l'erreur sur la part des 60 ans et plus. Une
rétention des seniors prolongée en tendance dégrade la prévision ($+0{,}33$\mesure{} point). La borne de
rétention $r \le 1{,}2$ d'un modèle voisin produit une croissance composée (+23\,\%\mesure{} de dentistes
en 2050 par rapport à la référence). Limite principale : le gain est petit, l'intervalle ignore la
corrélation entre origines, et les cibles ont déjà été vues lors d'un banc antérieur. L'article est une
première étape ; les familles de méthodes hors stock-flux restent à mesurer.
\end{abstract}

{\renewcommand\abstractname{Abstract}
\begin{abstract}
We look for an interpretable method that forecasts the private-sector (\emph{libéral}) headcount of
midwives, dentists, pharmacists, physiotherapists and nurses by French département better than a
stock-flow model with local rates. Five variants were registered before any run on a pre-registered
forecasting bench (rolling origins 2016--2025, mean absolute percentage error by département at five
years). All five are stock-flow models: same equation, different estimation. Hierarchical shrinkage of
the rates (département, region, nation) reaches 8.57\,\% MAPE against 8.68\,\% for the reference
($-0.11$ point, 95\,\% CI $[-0.21, -0.02]$), and the gain replicates on a second perimeter ($-0.10$).
Single-year ages tie on the total headcount; on simulated worlds they divide the error on the share
aged 60+ by three. Extrapolating the trend in senior retention degrades forecasts ($+0.33$ point). A
retention cap $r \le 1.2$ used by a related model yields compounding growth (+23\,\% dentists in 2050
relative to the reference). Main limits: the gain is small, the interval ignores correlation between
origins, and targets were already seen by an earlier bench. This is a first step; non stock-flow
families remain to be measured.
\end{abstract}}
\noindent\textbf{Mots-clés :} démographie des professions de santé ; projection ; stock-flux ;
shrinkage hiérarchique ; backtest pré-enregistré ; jumeau numérique.

\medskip
\noindent\textit{Provenance des nombres :} \mesure{} mesuré, \derive{} dérivé, \suppose{} supposé.
Tous viennent du rapport \texttt{v2/docs/lab/reports/2026-11-14-axe2-methodes.md}.

\section{Introduction}
Un décideur public qui anticipe des tensions territoriales a besoin d'une projection d'effectifs par
profession et par territoire, et de savoir ce qu'un levier (quotas de formation, âge de départ,
régulation de l'installation) y change. Le projet construit pour cela un \enquote{jumeau numérique} de
l'offre de soins : un état $N[\text{profession}, \text{secteur}, \text{territoire}, \text{sexe}, \text{âge}]$
que des modules de flux font avancer d'année en année. Son modèle de départ est un stock-flux par
composantes, comme les projections de la DREES~\cite{drees2021dd76}.

La question de cet article : une autre méthode, aussi lisible, prévoit-elle mieux ? \enquote{Lisible}
est une contrainte dure : chaque méthode doit nommer le paramètre qu'un levier modifie ; une méthode qui
ne le peut pas est exclue.

\paragraph{Ce que l'article apporte, et ce qu'il n'apporte pas.}
\begin{itemize}
  \item Le premier passage confirmatoire du protocole de prévision pré-enregistré du projet : le
        stock-flux à taux locaux bat la meilleure référence naïve (la tendance linéaire) de
        $-0{,}63$\mesure{} point de MAPE (IC 95\,\% $[-0{,}97\,;\,-0{,}27]$\mesure).
  \item Cinq variantes pré-inscrites, avec un verdict chacune. \textbf{Toutes sont des stock-flux} :
        elles gardent l'équation $N(t+1) = s \cdot A(N(t)) + I$ et ne changent que l'opérateur de
        vieillissement $A$, le partage d'information entre territoires ou la dépendance au temps de $s$.
        Les écarts mesurés (de l'ordre de 0,1\mesure{} point) en sont la conséquence.
  \item Un test sur mondes simulés à âges connus, une vérification de la borne de rétention d'un modèle
        voisin, et un diagnostic sur la répartition régionale des entrées.
  \item Les familles de méthodes hors stock-flux (microsimulation, bayésien hiérarchique ou espace-état,
        choix discret de l'installation) \emph{ne sont pas mesurées ici}. La section~\ref{sec:conclusion}
        en donne la feuille de route.
\end{itemize}

\section{Travaux existants}
\paragraph{Les modèles de projection d'effectifs de santé.} Une revue systématique de la littérature
2010--2023 retient 40 études qui projettent à la fois l'offre et la demande de
professionnels~\cite{lee2024hrh}. Les points qu'elle établit et qui nous concernent :
\begin{itemize}
  \item Côté offre, toutes recourent au stock-flux (40 sur 40) ; la simulation par agents (3 études) et
        la dynamique des systèmes (7) sont présentées comme des réponses à ses limites, notamment
        l'hétérogénéité de la population active, par exemple par âge.
  \item Seules 10 études travaillent à une échelle infranationale, et le manque de données locales y
        est un obstacle récurrent.
  \item La validation reste rare : 8 études valident la précision des prédictions contre des données
        passées, 13 font une analyse de sensibilité. Une erreur quadratique moyenne en pourcentage
        inférieure à 10\,\% y est donnée comme règle empirique d'acceptabilité.
  \item Le manque de standardisation dans la description des méthodes fait percevoir ces modèles
        comme des boîtes noires ; les auteurs proposent une grille de 30 items (cadre analytique,
        paramètres et leur estimation, hypothèses, incertitude, méthodes de validation, sources de
        données). Cet article en suit l'esprit : chaque paramètre, hypothèse et validation est
        explicite.
\end{itemize}
En France, les projections de la DREES pour les chirurgiens-dentistes, sages-femmes et pharmaciens
jusqu'en 2050 reposent sur un modèle par composantes (âge, sexe, mode d'exercice) : cessations,
transitions de mode d'exercice, entrées décalées de la durée des études, diplômés à
l'étranger~\cite{drees2021dd76}. Notre référence en reprend la structure, simplifiée à une rétention
nette et à un flux d'entrées nettes par cellule.

\paragraph{Cohortes par composantes en bayésien, à l'échelle infranationale.} Alexander et Alkema
placent l'identité des cohortes par composantes dans un cadre bayésien pour estimer et projeter des
populations de comtés kényans avec très peu de données~\cite{alexander2021bayesian}. Les processus
(mortalité par composantes principales des tables nationales, migration nette à profil d'âge constant
et volume en marche aléatoire, multiplicateur âge-temps résiduel) sont reliés aux observations par
des modèles de données qui portent l'erreur d'échantillonnage, et la somme des comtés est contrainte
autour du total national publié. Validées hors échantillon contre le recensement 2019, ses
projections sont à environ 1\,\% des effectifs observés, avec une erreur quadratique 3 à 10 fois
plus faible qu'une extrapolation linéaire, et une couverture des intervalles à 90\,\% proche de sa
valeur nominale (cité). C'est la forme la plus proche de notre jumeau parmi les méthodes hors
stock-flux déterministe : même identité comptable, mais des taux estimés avec un modèle explicite,
une incertitude propagée et un partage d'information entre zones. Notre shrinkage A2 en est la
version déterministe et sans incertitude.

\paragraph{Ce dont nous partons.} Un observatoire construit par un membre de
l'équipe~\cite{evolutionsoin2026} utilise un stock-flux voisin, avec un shrinkage département $\to$
région $\to$ national (poids 200) et des taux de maintien bornés à $[0{,}5\,;\,1{,}2]$. Son audit
signale que le vieillissement uniforme dans la tranche biaise la part des seniors de +8 à +19\,\% sur
des mondes simulés (chiffre cité). Nous testons ces trois éléments.

\paragraph{Ce que nous ne citons pas encore.} Un état de l'art interne au
projet\footnote{\texttt{v2/docs/lab/explorations/2026-11-14-etat-de-lart-methodes.md}.} a repéré
d'autres familles en boîte blanche : modèles âge-période-cohorte bayésiens pour les cessations,
logit conditionnel du lieu d'installation, réconciliation optimale des prévisions entre échelles,
règles de score propres et histogrammes PIT pour valider des prévisions probabilistes. Ces références
n'ont pas encore été lues en texte intégral ; la règle du projet les exclut de la bibliographie tant
que ce n'est pas fait.

\section{Données}
\begin{itemize}
  \item \textbf{Effectifs.} DREES, démographie des professionnels de santé depuis 2012, publication du
        2026-07-02 (répertoire RPPS) : profession $\times$ département $\times$ sexe $\times$ tranche
        d'âge quinquennale $\times$ mode d'exercice $\times$ secteur d'activité, 2012--2026. Kinés : RPPS
        à partir de 2017 seulement (rupture ADELI $\to$ RPPS). Infirmiers libéraux : 2012--2023.
  \item \textbf{Séparation public / privé.} Construite par l'axe 4 du projet à partir du secteur
        d'activité des fichiers DREES ; les salariés ambigus sont répartis par leur part privée
        observée dans FINESS (extraction du 2026-05-04). Nous utilisons la variante centrale.
  \item \textbf{Deux périmètres.} \emph{Privé} (confirmatoire) : secteur privé, lucratif ou non ; pour
        les infirmiers, l'âge des salariés privés manque et seuls les infirmiers libéraux sont gardés.
        \emph{Ancien} (secondaire) : modes libéral et mixte, pour la comparaison avec un banc antérieur.
  \item \textbf{Inscrits en formation.} DREES, fichier \texttt{990\_drees\_mmop2\_f.xlsx}, feuille
        \texttt{Feuil1} (2020--2024), qui porte la correspondance université $\to$ ARS.
\end{itemize}
Départements : 101, hors collectivités d'outre-mer (codes 999 et 20 exclus).

\section{Méthode}
\subsection{Référence (A0)}
Par cellule $c$ = (profession, département, sexe) et tranche $a$ :
\begin{equation}
  N_{c}(t+1) = s_{c} \odot A\big(N_{c}(t)\big) + I_{c},
\end{equation}
où $A$ fait passer $1/5$ de chaque tranche à la suivante (âges supposés uniformes dans la tranche),
$s_{c,a} = \min(1, r_{c,a})$ est la rétention nette ajustée par moindres carrés sur la fenêtre, et
$I_{c,a}$ est le flux net moyen là où le ratio brut $r_{c,a}$ dépasse~1. Les taux départementaux sont
rétrécis vers le national : $r = (\text{num} + k\,r_{\text{nat}})/(\text{den} + k)$, $k = 200$\suppose.

\subsection{Variantes inscrites}
\begin{description}
  \item[A1a, âges interpolés.] Chaque année, chaque tranche est désagrégée en âges simples (densité
        linéaire par morceaux entre milieux de tranches, masse de chaque tranche conservée), vieillie
        d'un an, puis réagrégée. Le même opérateur sert à l'ajustement et à la projection. Hypothèses :
        tranche \enquote{moins de 25 ans} = 20--24 ans\suppose, tranche ouverte 65+ de largeur
        10 ans\suppose.
  \item[A1b, âges simples.] L'état est tenu en âges simples (20 à 74 ans et plus), désagrégé une fois à
        l'origine ; les cohortes sont suivies exactement. Rétention et entrées restent calibrées par
        tranche ; les entrées sont réparties uniformément dans la tranche\suppose.
  \item[A2, shrinkage hiérarchique.] $r_{\text{dép}} = (\text{num} + k\,r_{\text{rég}})/(\text{den} + k)$
        et $r_{\text{rég}} = (\text{num}_{\text{rég}} + k\,r_{\text{nat}})/(\text{den}_{\text{rég}} + k)$,
        $k = 200$\suppose{} aux deux niveaux, non réglé sur le banc.
  \item[A3, rétention des 60 ans et plus en tendance.]
        $s_{c,a}(t) = \min\!\big(1,\ s_{c,a} + \lambda\, b_{p,\text{sexe},a}\,(t - \bar t)\big)$ pour
        $a \ge 60$, où $b$ est la pente nationale des ratios annuels sur la fenêtre et $\bar t$ son
        milieu ; $\lambda = 1$.
  \item[A123.] A1a + A2 + A3, combinaison pré-spécifiée.
  \item[B4 (diagnostic, non candidat).] Schéma de l'observatoire~\cite{evolutionsoin2026} : $r = 1$
        et entrées = résidus moyens sur 5 ans dans les tranches de moins de 35 ans ; ailleurs $r$
        rétréci département $\to$ région $\to$ national et borné à $[0{,}5\,;\,1{,}2]$, sans entrée
        absolue.
\end{description}

\subsection{Ce que chaque levier modifie}
\begin{table}[h]
\centering\small
\begin{tabular}{@{}p{0.2\textwidth}p{0.45\textwidth}p{0.27\textwidth}@{}}
\toprule
Variante & Paramètre que touche un levier & Nature \\
\midrule
A0, A1a, A1b, A2 & \texttt{methodes\_sorties.}\allowbreak\texttt{decalage\_retraite}, \texttt{cumul\_multiplicateur} & rétention $s$ par tranche \\
 & \texttt{methodes\_entrees.multiplicateur} & flux d'entrées $I$ \\
A3 & \texttt{methodes\_sorties.tendance\_60plus} ($\lambda$) & pente de rétention des 60+ \\
A2 & $k$ (poids du shrinkage) & réglage d'estimation, pas un levier \\
\bottomrule
\end{tabular}
\caption{Paramètres touchés par les leviers. Aucune variante n'introduit de paramètre sans nom.}
\end{table}

\subsection{Mise en \oe{}uvre}
Les variantes sont trois modules (\texttt{methodes\_vieillissement}, \texttt{methodes\_sorties},
\texttt{methodes\_entrees}) branchés sur le registre du jumeau, sans modifier son code. Aux options de
la référence, la composition reproduit le jumeau de référence à $9 \times 10^{-13}$\mesure{} personne
près ; la désagrégation conserve la masse à $5 \times 10^{-13}$\mesure{} près. Le recalcul d'un
scénario jusqu'en 2050 (5 professions, département) prend de 14 à 40\mesure{}~ms.

\section{Résultats}
\subsection{Banc de prévision, périmètre privé (confirmatoire)}
\begin{table}[h]
\centering\small
\resizebox{\textwidth}{!}{%
\begin{tabular}{@{}lrrr@{}}
\toprule
Variante & MAPE $h=5$ (\%) & $\Delta$ vs tendance (IC 95\,\%) & $\Delta$ vs A0 (IC 95\,\%) \\
\midrule
A0 stock-flux taux locaux & 8,68 & $-0{,}63$ $[-0{,}97\,;\,-0{,}27]$ & --- \\
\textbf{A2 shrinkage hiérarchique} & \textbf{8,57} & $-0{,}74$ $[-1{,}07\,;\,-0{,}40]$ & $\mathbf{-0{,}11}$ $[-0{,}21\,;\,-0{,}02]$ \\
A1b âges simples & 8,70 & $-0{,}60$ $[-0{,}96\,;\,-0{,}24]$ & $+0{,}03$ $[-0{,}07\,;\,+0{,}11]$ \\
A123 combinaison & 8,74 & $-0{,}56$ $[-0{,}92\,;\,-0{,}21]$ & $+0{,}06$ $[-0{,}07\,;\,+0{,}18]$ \\
A1a âges interpolés & 8,77 & $-0{,}53$ $[-0{,}89\,;\,-0{,}17]$ & $+0{,}09$ $[+0{,}01\,;\,+0{,}18]$ \\
A3 rétention 60+ en tendance & 9,00 & $-0{,}30$ $[-0{,}67\,;\,+0{,}08]$ & $+0{,}33$ $[+0{,}23\,;\,+0{,}42]$ \\
\textit{B4 observatoire (diagnostic)} & \textit{8,32} & $-0{,}98$ $[-1{,}37\,;\,-0{,}59]$ & $-0{,}35$ $[-0{,}65\,;\,-0{,}08]$ \\
tendance linéaire (naïve) & 9,30 & --- & --- \\
persistance (naïve) & 9,92 & --- & --- \\
\bottomrule
\end{tabular}}
\caption{MAPE par département à cinq ans, départements de 30 professionnels ou plus à l'origine,
moyenne à poids égal sur 22 cellules profession $\times$ origine. Toutes les valeurs sont
mesurées\mesure. IC par bootstrap des départements (2\,000 tirages, graine 20261114).}
\label{tab:prive}
\end{table}

Toutes les variantes, la référence comprise, battent la tendance, sauf A3 (égalité). Les victoires
tiennent aussi à l'IC 99\,\% (correction de Bonferroni pour cinq variantes). La limite $h \to 0$ est
respectée : à $h = 1$, aucune variante n'est à plus d'un point de la persistance.

\paragraph{Par profession.} Le gain de A2 vient surtout des dentistes (8,66 contre 8,95\mesure) et des
infirmiers libéraux (5,72 contre 5,89\mesure) ; il est nul ou négatif pour les kinés et les pharmaciens.
Les sages-femmes privées ont la MAPE la plus forte (13 à 14\,\%\mesure), dans une série libérale qui a
été multipliée par 1,8\mesure{} entre 2012 et 2026. Pour les dentistes, la référence ne bat la tendance
que sur 2 origines sur 6\mesure : leur classement dépend de l'origine.

\paragraph{Prédictions écrites avant le banc.} Trois prédictions avaient été déposées
(empreinte SHA-256 de la dérivation \texttt{1c324f5795b0}) :
$|\Delta(\text{A1}, \text{A0})| < 0{,}3$ point (observé $+0{,}09$ et $+0{,}03$\mesure : tient) ;
A2 améliore A0 de 0 à 0,3 point (observé $-0{,}11$\mesure : tient) ;
A3 réduit le biais national des pharmaciens et des dentistes (pharmaciens : $-0{,}15$ $\to$
$+2{,}28$\,\%\mesure{} : \textbf{falsifiée}).

\subsection{Ancien périmètre (secondaire)}
Sur le périmètre libéral + mixte, A2 reste le seul gain : $-0{,}10$\mesure{} point
(IC 95\,\% $[-0{,}21\,;\,-0{,}002]$\mesure). A1b est à égalité ($-0{,}03$\mesure), A3 et la combinaison
sont pires ($+0{,}76$ et $+0{,}50$\mesure). \textbf{B4 se renverse} : meilleur que A0 au privé, il est
pire de $+0{,}53$\mesure{} point ici. La référence y mesure 8,21\,\%\mesure, contre 8,08\,\% dans le banc
antérieur du projet, sur une grille différente.

\subsection{Mondes simulés : le biais du vieillissement uniforme}
Dix-huit mondes à âges simples connus (un pic de génération de taille 0, 60 ou 120, centré à 40, 50 ou
55 ans ; entrées constantes ou en hausse de 2\,\% par an\suppose) sont observés en tranches, ajustés
sur dix ans, puis projetés.
\begin{table}[h]
\centering\small
\resizebox{\textwidth}{!}{%
\begin{tabular}{@{}lrrr@{}}
\toprule
Méthode & $|$erreur$|$ 60+ à $h=5$, moy. (max) & $|$erreur$|$ 60+ à $h=10$ & $|$erreur$|$ total à $h=5$ \\
\midrule
uniforme (A0) & 7,0 (18,9) & 5,2 (12,0) & 3,9 \\
A1a âges interpolés & 5,3 (13,5) & 4,5 (9,7) & 3,0 \\
A1b âges simples & 2,1 (5,9) & 1,8 (2,9) & 3,8 \\
\bottomrule
\end{tabular}}
\caption{Erreurs en \%, mesurées\mesure{} sur mondes simulés.}
\end{table}
Le vieillissement uniforme biaise la part des 60 ans et plus jusqu'à 19\,\%\mesure, l'ordre de grandeur
signalé par l'observatoire. Le biais ne se voit pas sur le total à cinq ans : la rétention ajustée
l'absorbe dans la fenêtre\derive, et l'erreur totale est dominée par les entrées. C'est pourquoi A1b
fait jeu égal au banc tout en étant plus juste sur la structure par âge.

\subsection{La borne $r \le 1{,}2$ fait-elle s'emballer les entrées ?}
De 27 à 47\,\%\mesure{} des cellules de 35 ans et plus ont $r > 1$ sous B4 : leur croissance est
proportionnelle au stock et se compose. De 2026 à 2050, l'écart de croissance annuelle à A0 va de
$+0{,}1$ à $+0{,}9$\mesure{} point ; pour les dentistes, l'effectif 2050 vaut 1,61 fois\mesure{}
celui de 2026 contre 1,31\mesure{} sous A0 (+23\,\%\mesure). Ce n'est pas l'emballement de
+3,5\,\% par an observé lors du banc antérieur, mais le piège est en partie reproduit, et il se
concentre au-delà de l'horizon que l'on peut tester rétrospectivement. Une ablation exploratoire
(conçue après le banc) attribue l'avantage de B4 au privé à cette croissance proportionnelle, non aux
entrées récentes ni à la hiérarchie.

\subsection{La région de formation prédit-elle la région d'installation ?}
Sur 12 régions, la part du stock 2026 prédit la part régionale des entrées nettes avec une
corrélation de 0,996\mesure{} (dentistes et pharmaciens), contre 0,64\mesure{} (dentistes) et
0,94\mesure{} (pharmaciens) pour la part des étudiants (années 2 à 6, 2020--2024). En niveau, répartir
les entrées par région de formation dégraderait donc la prévision. Le test des \emph{variations}
demande des inscrits par université avant 2020, que nous n'avons pas.

\section{Validation et menaces à la validité}
\paragraph{Protocole.} Le banc suit le protocole de prévision du projet, gelé par empreinte avant tout
calcul : origines glissantes, horizons 1, 3, 5, 8 et 10 ans, métrique principale unique (MAPE par
département à cinq ans), critère = borne haute de l'IC 95\,\% de l'écart à la meilleure naïve
inférieure à 0. Les cinq variantes ont été inscrites au registre avant leur premier passage ; B4 et
l'ablation sont déclarés diagnostic et exploratoire. La comparaison à A0 n'est pas le critère du
protocole : elle était annoncée comme exploratoire dans la dérivation.

\paragraph{Menaces.}
\begin{itemize}
  \item \textbf{Petite marge.} Le gain de A2 ($-0{,}11$ point) a une borne haute de $-0{,}02$ (privé)
        et $-0{,}002$ (ancien). Le bootstrap par département ignore la corrélation entre origines qui
        se recouvrent : l'intervalle est probablement trop étroit\suppose.
  \item \textbf{Contamination.} Les cibles 2023 et 2026 ont déjà été vues par un banc antérieur du
        projet. Seul un dépôt de prévisions 2027--2030, scoré à la publication du RPPS 2027, sera hors
        échantillon.
  \item \textbf{Périmètre privé.} La répartition public / privé des salariés ambigus repose sur une
        part FINESS de 2026 appliquée à toute la série ; seule la variante centrale a été passée.
  \item \textbf{Hypothèses d'âge.} Les largeurs des tranches extrêmes et la répartition uniforme des
        entrées dans la tranche sont supposées.
  \item \textbf{Leviers.} Une revue du jumeau a relevé un biais du levier retraite sur la tranche 65+,
        hérité par \texttt{methodes\_sorties}. Aucun levier n'a servi à départager les méthodes.
\end{itemize}

\section{Reproductibilité}
Depuis la racine du dépôt :
\begin{verbatim}
export PYTHONPATH=v2
uv run python -m methodes.verif     # contrôles V1-V5
uv run python -m methodes.banc      # banc N1
uv run python -m methodes.monde     # mondes simulés
uv run python -m methodes.b4        # borne r <= 1,2
uv run python -m methodes.c5        # région de formation
uv run python -m methodes.ablation  # exploratoire
\end{verbatim}
Graine du bootstrap 20261114, 2\,000 tirages. Sorties dans \texttt{v2/methodes/out/}. Données :
\texttt{data/processed/stock.parquet} (DREES, 2026-07-02) et
\texttt{v2/data/processed/secteur\_stock.parquet} (reconstruit par \texttt{v2/data/build\_secteur.py}).
Le banc de chaque périmètre prend environ 20\mesure{}~s, plus quelques minutes de bootstrap. Le code
n'est pas encore versionné (aucun commit au 2026-11-14).

\section{Conclusion}
\label{sec:conclusion}
Sur le banc pré-enregistré, le stock-flux à taux locaux bat les références naïves, et un shrinkage
hiérarchique l'améliore de peu mais de façon reproductible sur deux périmètres. Les âges fins n'aident
pas la prévision de l'effectif total ; ils corrigent la structure par âge, dont dépendent les leviers
retraite. Prolonger la tendance de rétention des seniors dégrade la prévision.

Ces résultats restent internes à une seule famille : \textbf{les cinq variantes sont des stock-flux}.
Leur proximité (écarts de l'ordre de 0,1 point) dit que l'estimation des sorties n'est plus le
facteur limitant. Feuille de route vers d'autres familles, par ordre de priorité :
\begin{enumerate}
  \item \textbf{Choix discret de l'installation} (logit sur des attributs départementaux datés :
        croissance de la population, densité relative, zonage). C'est la famille qui vise la source
        d'erreur dominante, la répartition des entrées, et son coefficient de zonage est le levier
        régulation. Il faut une série de population par département avant 2018 (\textasciitilde{}1\,h).
  \item \textbf{Bayésien hiérarchique ou espace-état}, évalué à la commune selon le protocole
        d'échelle fine, où les petits effectifs rendent le partage d'information décisif
        (\textasciitilde{}2\,h).
  \item \textbf{Microsimulation individuelle}, pour les bandes d'incertitude et les règles non
        linéaires (régulation par commune) ; sa prévision ponctuelle égale celle du stock-flux à taux
        égaux\derive.
\end{enumerate}
Autres suites : inscrire et passer la combinaison A2 + A1b (\textasciitilde{}10\,min), déposer les
prévisions 2027--2030 (\textasciitilde{}10\,min), et trancher entre A0 et B4 à 2050 par la
reproduction des projections DREES~\cite{drees2021dd76} (\textasciitilde{}30\,min).

\bibliographystyle{plainnat}
\bibliography{refs}
\appendix
\section{Annexes}
\paragraph{Désagrégation.} Pour des tranches de bornes basses $\ell_a$ et de largeurs $w_a$, la densité
moyenne $m_a = N_a / w_a$ est placée au milieu de la tranche ; la densité de l'âge simple $x$ est
l'interpolation linéaire de $m$ en $x + 1/2$, tronquée à 0, puis remise à l'échelle pour que chaque
tranche garde sa masse $N_a$. Avec une densité uniforme dans la tranche, l'opérateur $A$ redonne
exactement $0{,}8\,N_a + 0{,}2\,N_{a-1}$ (vérifié à $1{,}6 \times 10^{-14}$\mesure{} près).
\end{document}
